\clearpage\begingroup
\providecommand{\ObsSec}[2]{#1\mathord{\overset{\mathrm{s}}{,}}#2}
\providecommand{\ObsArc}[2]{#1\mathord{\overset{\prime\prime}{,}}#2}
\providecommand{\ObsJupiter}{\mathord{\tikz[baseline=-.1ex,x=1em,y=1em,line width=.045em]{
  \draw (.02,.48) .. controls (.04,.70) and (.33,.69) .. (.31,.46) .. controls (.29,.29) and (.14,.19) .. (.02,.13)--(.49,.13);
  \draw (.37,.64)--(.37,-.06);}}}
\providecommand{\ObsVesta}{\mathord{\tikz[baseline=-.1ex,x=1em,y=1em,line width=.045em]{
  \draw (.04,.02) rectangle (.40,.38);
  \draw (.22,.38)--(.22,.48) (.12,.63)--(.22,.48)--(.32,.63);}}}

\newfontfamily\ReichenbachGreek{STIXTwoMath-Regular.otf}
\newcommand{\ReichenbachBeta}{\text{\ReichenbachGreek ϐ}}
\SourceLeaf{287}\NoSourceHead
\begin{center}\large
BEOBACHTUNGEN MIT\\
EINEM 12-ZÖLLIGEN REICHENBACHSCHEN\\
KREISE ZUR BESTIMMUNG DER POLHÖHE\\
DER [ALTEN] GÖTTINGER STERNWARTE.
\end{center}
\begin{center}\small
\rule{.78\textwidth}{.35pt}\\[.2em]
Monatliche Correspondenz, Band XXVII, S. 481—484, 1813 Mai.\\[-.2em]
\rule{.78\textwidth}{.35pt}
\end{center}
\begin{center}\textls[120]{Aus einem Schreiben des Hrn. Prof. \textsc{Gauss}.}\quad Göttingen, 18. Mai 1813.\end{center}
\noindent\hspace*{\parindent}Sie wissen, theuerster Freund, welche zufällige Hindernisse den Anfang der astronomischen Beobachtungen mit dem herrlichen \textsc{Reichenbach}schen Kreise, welchen unsere Sternwarte gegen Ende des vorigen Jahres erhielt, bis zur Mitte des März verzögert haben. Ich habe jetzt das Vergnügen, Ihnen die Erstlinge meiner Beobachtungen mitzutheilen, welche die Bestimmung der Polhöhe der Sternwarte zum Zweck haben. Auf der Südseite des Meridians, wo ich dem Kreise eine eben so feste Aufstellung wie er jetzt für die Nordseite hat, noch nicht habe geben können, sind noch keine Beobachtungen angestellt.
\begin{center}
Zenith-Distanzen des Polarsterns in der untern Culmination.
\medskip
\begingroup\fontsize{9.5}{12}\selectfont\setlength{\tabcolsep}{4pt}\renewcommand{\arraystretch}{1.25}
\begin{tabular}{r|c|c|c|c|c}
\hline\hline
1813&\begin{tabular}{@{}c@{}}Scheinb. Zenith-Distanz\\in der untern Culmin.\end{tabular}&\begin{tabular}{@{}c@{}}Anz. der\\Beob.\end{tabular}&Barometer&\begin{tabular}{@{}c@{}}Thermo-\\meter\end{tabular}&\begin{tabular}{@{}c@{}}Re-\\fraction\end{tabular}\\
\hline\hline
März 20.&$40^0\,8'\,\ObsArc{35}{64}$&10&$27\,\mathrm Z\quad8\mathord{\overset{\mathrm l}{,}}9$&$+2\mathord{\overset{0}{,}}5$&$\ObsArc{49}{84}$\\
22.&$35{,}14$&18&$27\qquad8{,}3$&$+1{,}25$&$50{,}05$\\
26.&$33{,}67$&18&$28\qquad1{,}2$&$+0{,}5$&$51{,}01$\\
31.&$38{,}60$&18&$27\qquad5{,}7$&$+5{,}5$&$48{,}70$\\
April 3.&$38{,}05$&18&$27\qquad7{,}8$&$+1{,}75$&$49{,}86$\\
7.&$41{,}05$&32&$27\qquad10{,}3$&$+5{,}5$&$49{,}39$\\
8.&$39{,}94$&22&$27\qquad9{,}9$&$+8{,}25$&$48{,}65$
\end{tabular}\endgroup
\end{center}
\clearpage\SourceLeaf{288}\HistoricalHead{283}{R}{BEOBACHTUNGEN MIT EINEM 12-ZÖLLIGEN REICHENBACHSCHEN KREISE.}
\noindent Hieraus also:
\begin{center}\begingroup\fontsize{9.5}{12}\selectfont\setlength{\tabcolsep}{4pt}\renewcommand{\arraystretch}{1.25}
\begin{tabular}{r|c|c|c}
\hline\hline
1813&\begin{tabular}{@{}c@{}}Wahre Zenith-Distanz\\des Polarsterns\end{tabular}&\begin{tabular}{@{}c@{}}Scheinbare\textsuperscript{[*)]}\\Polar-Distanz\end{tabular}&Polhöhe\textsuperscript{[*)]}\\
\hline\hline
März 20.&$40^0\,9'\,\ObsArc{25}{48}$&$1^0\,41'\,\ObsArc{20}{34}\;(\ObsArc{19}{63})$&$51^0\,31'\,\ObsArc{54}{86}\;(\ObsArc{54}{15})$\\
22.&$25{,}19$&$20{,}92\;(20{,}24)$&$55{,}73\;(55{,}05)$\\
26.&$24{,}68$&$22{,}09\;(21{,}48)$&$57{,}41\;(56{,}80)$\\
31.&$27{,}33$&$23{,}58\;(23{,}02)$&$56{,}25\;(55{,}69)$\\
April 3.&$27{,}91$&$24{,}46\;(23{,}93)$&$56{,}55\;(56{,}02)$\\
7.&$30{,}44$&$25{,}64\;(25{,}16)$&$55{,}20\;(54{,}72)$\\
8.&$28{,}59$&$25{,}92\;(25{,}46)$&$57{,}33\;(56{,}87)$\\
\multicolumn{3}{l}{Mittel aus 136 Beobachtungen}&$51^0\,31'\,\ObsArc{56}{20}\;(\ObsArc{55}{64})$.
\end{tabular}\endgroup
\end{center}
\noindent\hspace*{\parindent}Bei Berechnung dieser Beobachtungen liegt die Polar-Distanz $\alpha$ im kleinen Bär für den Anfang von 1812
\[=1^0\,41'\,\ObsArc{41}{74}\]
\noindent zum Grunde\textsuperscript{**)}. Herr \textsc{Pond} hat sie mit dem neuen \textsc{Troughton}schen Kreise in Greenwich für dieselbe Epoche
\begin{center}\begin{tabular}{ll}
aus Sommerbeobachtungen&$=1^0\,41'\,\ObsArc{41}{60}$\\
aus Winterbeobachtungen&$=1\quad41\quad41{,}00$
\end{tabular}\end{center}
\noindent gefunden; die Pariser Astronomen haben mit dem \textsc{Reichenbach}schen dreifüssigen Kreise erhalten:
\[1^0\,41'\,\ObsArc{41}{19}.\]
\noindent\hspace*{\parindent}Endlich folgt aus den Rechnungs-Resultaten, die Herr von \textsc{Zach} aus \textsc{Orianis} Beobachtungen gezogen, und Mon. Corresp. Bd. XXVII, S. 110 mitgetheilt hat, eine Verminderung der Polar-Distanz gegen von \textsc{Zachs} eigene Bestimmung von $\ObsArc{0}{77}$, also für $1812=1^0\,41'\,\ObsArc{40}{97}$. Das Mittel wäre also aus von \textsc{Zachs}, \textsc{Ponds}, \textsc{Orianis} und den Bestimmungen der französischen Astronomen
\[1^0\,41'\,\ObsArc{41}{30}\]
\noindent genau mit dem Mittel aus \textsc{Ponds} Resultate harmonirend, wonach obige Polhöhe um $\ObsArc{0}{44}$ zu vermindern wäre. Die Reduction auf den Mittelpunkt der Sternwarte beträgt noch $-\ObsArc{0}{16}$; wir haben also die Polhöhe der Göttinger
\begingroup\renewcommand{\thefootnote}{[*)}\footnotetext{Die eingeklammerten Zahlen hat \textsc{Gauss} in seinem Handexemplar handschriftlich eingetragen.]}\endgroup
\begingroup\renewcommand{\thefootnote}{**)}\footnotetext{Wie ich sie aus Combination von Herrn von \textsc{Zachs} neuester Angabe für 1810 mit den Bestimmungen von \textsc{Méchain} und \textsc{Delambre} angenommen hatte.}\endgroup
\HistoricalFoot{}{$36^{*}$}
\clearpage\SourceLeaf{289}\HistoricalHead{284}{L}{PRAKTISCHE UND SPHÄRISCHE ASTRONOMIE. VERÖFFENTLICHUNGEN UND BRIEFWECHSEL.}
\noindent Sternwarte $=51^0\,31'\,\ObsArc{56}{04}$ oder $=51^0\,31'\,\ObsArc{55}{60}$ $(51^0\,31'\,\ObsArc{55}{48})$ je nachdem wir für die Declination des Polarsterns blos Herrn von \textsc{Zachs} Bestimmung, oder das Mittel von vier Astronomen annehmen. \textsc{Mayers} Bestimmung dieser Polhöhe wäre demnach nur $2''$ oder $\ObsArc{1}{6}$ zu klein.

Ich habe in diese Reihe die Beobachtungen von drei andern Abenden nicht mit aufgenommen, weil dies meine \textls[120]{allerersten} Versuche waren; inzwischen stimmen auch diese auf das schönste überein, es sind folgende:
\begin{center}\begingroup\small\setlength{\tabcolsep}{3pt}
\begin{tabular}{rrrlr}
1813&März&14.&Vier Zenith-Distanzen des Polarsterns \textls[120]{ausser}&\\
&&&\quad der Culmination geben Polhöhe\quad\mbox{. . . . . .}&$51^0\,31'\,\ObsArc{56}{57}$\\[.3em]
»&»&15.&Zenith-Distanz von $\ReichenbachBeta$ Cephei in der untern&\\
&&&\quad Culmination 4 Beobachtungen.\,\mbox{. . . . . . .}&$56{,}14$\\[.3em]
»&»&17.&Zenith-Distanz von $\ReichenbachBeta$ Cephei in der untern&\\
&&&\quad Culmination 6 Beobachtungen.\,\mbox{. . . . . . .}&$54{,}77$\\[-.4em]
\multicolumn{5}{r}{\rule{1.35in}{.35pt}}\\[-.2em]
\multicolumn{4}{r}{Mittel aus 14 Beobachtungen}&$51^0\,31'\,\ObsArc{55}{68}$.
\end{tabular}\endgroup\end{center}
\noindent\hspace*{\parindent}Die Beobachtungen vom 14. waren bei solchen Stundenwinkeln angestellt, wo die Declination wenig Einfluss auf die Rechnung hatte; über die Art der Berechnung erkläre ich mich bei einer andern Gelegenheit umständlicher. Der Polarstern ist ein am Firmament so wohl bestimmter Punkt, dass man ihn zu jeder Stunde mit gleicher Schärfe zu Breitenbestimmungen anwenden kann; ja ich glaube, dass es in Rücksicht auf \textls[120]{das Beobachten} noch vortheilhafter wäre, ihn immer in beträchtlicher Entfernung vom Meridian zu beobachten. Bekanntlich stören der Beobachter und der, welcher das Niveau einstellt, einander immer wechselseitig etwas, und zwar mehr, wenn ein Object von ganz oder fast ganz constanter Höhe observirt wird, da hingegen, wenn die Höhe im Zunehmen oder Abnehmen begriffen ist, der Beobachter, sobald das Niveau gut steht, die \textls[120]{schärfste} Bissection durch \textls[120]{Abwarten} erreichen kann. Die Beobachtungen in der Culmination haben nur den Vorzug einer kürzern einfachern Berechnung und der Unabhängigkeit von der Rectascension des Sterns, letztere ist aber so gut bestimmt, dass der daraus entspringende Fehler als $0$ angesehen werden darf; man weicht ihm überdies auch aus, wenn man zwei \textls[120]{Sets} von Beobachtungen ungefähr in gleichen Abständen vor und nach der Culmination anstellt. Doch hiervon ein andermal ausführlicher. Die Declination von $\ReichenbachBeta$ Cephei ist nach \textsc{Piazzi} angenommen.
\begin{center}\rule{.31\textwidth}{.35pt}\end{center}
\clearpage\endgroup
